% Euler buckling — analytical derivation
% AIM Interactive Mechanics series. Compile with pdflatex.
\documentclass[10pt,a4paper]{article}
\usepackage[margin=2cm]{geometry}
\usepackage{amsmath}
\usepackage{booktabs}
\usepackage[T1]{fontenc}
\usepackage{lmodern}
\usepackage{microtype}
\usepackage{titlesec}
\titlespacing*{\section}{0pt}{9pt}{2pt}
\usepackage{xcolor}
\definecolor{iron}{HTML}{935636}
\definecolor{celdark}{HTML}{4B6C61}
\definecolor{celline}{HTML}{CDE0D8}
\usepackage[colorlinks=true,urlcolor=iron,linkcolor=black]{hyperref}
\setlength{\parindent}{0pt}
\setlength{\parskip}{4pt}
\AtBeginDocument{\setlength{\abovedisplayskip}{6pt}\setlength{\belowdisplayskip}{6pt}\setlength{\abovedisplayshortskip}{3pt}\setlength{\belowdisplayshortskip}{3pt}}
\pagestyle{empty}

\begin{document}

{\footnotesize\textcolor{celdark}{\textsc{aim interactive mechanics \,\textperiodcentered\, analytical derivation}}}

{\LARGE Euler Buckling}

\textcolor{celdark}{\small Companion to the interactive simulator at
\url{https://aimlab.uk/sims/buckling.html}}

\medskip
{\color{celline}\hrule height 1pt}

\section*{Setup and assumptions}
A straight prismatic member of length $L$ and rectangular cross-section
$w \times h$ carries an axial compressive force $P$; the material is linear
elastic with Young's modulus $E$. We address the question of ``at what load does the straight configuration stop being the only equilibrium?''. This is an \emph{eigenvalue} problem.

\section*{Weak axis}
The rectangle has two centroidal second moments,
$I_{xx} = w h^3/12$ and $I_{yy} = h w^3/12$, for bending in the $h$- and
$w$-directions respectively. A compressed member bows whichever way is easiest
--- about the weak axis,
$I_{\min} = \min(I_{xx}, I_{yy})$: a ruler squashed end-to-end always bows in
its thin direction.

\section*{Governing equation (pinned--pinned)}
Suppose the member has deflected laterally by a small $u(x)$. Moment equilibrium
of the segment beyond a cut at $x$, taken about the cut, shows the axial force now
produces an internal bending moment $M(x) = -P\,u(x)$. With moment--curvature
$M = EI\,u''$,
\begin{equation}
EI\,u'' + P\,u \;=\; 0
\qquad\Longrightarrow\qquad
u'' + k^2 u = 0, \quad k^2 = \frac{P}{EI},
\end{equation}
with $u(0) = u(L) = 0$ for pinned ends. The general solution is
$u = A\sin kx + B\cos kx$; the condition $u(0)=0$ gives $B = 0$, and $u(L)=0$
then requires $A \sin kL = 0$. Either $A = 0$ (the straight state --- always an
equilibrium) or $\sin kL = 0$, i.e.\ $kL = n\pi$. A bent equilibrium first
becomes possible at $n = 1$:
\begin{equation}
\boxed{\; P_{\mathrm{cr}} \;=\; \frac{\pi^2 E I_{\min}}{L^2} \;}
\qquad\text{with mode shape}\qquad
u(x) = A \sin\!\frac{\pi x}{L}.
\end{equation}
The amplitude $A$ is \emph{indeterminate}: linearised theory predicts the
critical load, not the actual displacement.

\section*{Other end conditions}
With general supports the governing form is fourth order,
$EI\,u'''' + P\,u'' = 0$, with solution $u = A\sin kx + B\cos kx + Cx + D$
fitted to the four end conditions. Each case reduces to the pinned--pinned
result with an \emph{effective length} $KL$ --- the distance between inflection
points of its mode: $P_{\mathrm{cr}} = \pi^2 E I_{\min} / (K L)^2$.
\begin{center}
\begin{tabular}{lll}
\toprule
End conditions & Characteristic equation & $K$ \\
\midrule
pinned--pinned & $\sin kL = 0$            & $1$ \\
fixed--free    & $\cos kL = 0$            & $2$ \\
fixed--pinned  & $\tan kL = kL$ \;($kL = 4.493$) & $\pi/4.493 = 0.699$ \\
fixed--fixed   & $\cos kL = 1$            & $1/2$ \\
\bottomrule
\end{tabular}
\end{center}
Clamping the ends quadruples the capacity; a flagpole geometry carries only a quarter.

\section*{Worked example (simulator defaults)}
$E = 3$~GPa (acrylic), $w = 8$~mm, $h = 25$~mm, $L = 500$~mm, pinned--pinned.
The weak axis is $y$--$y$: $I_{yy} = 25 \times 8^3 / 12 = 1067~\text{mm}^4$, well
below $I_{xx} = 10\,417~\text{mm}^4$. Then
\begin{equation}
P_{\mathrm{cr}} = \frac{\pi^2 \times 3\times10^{9} \times 1.067\times10^{-9}}{0.5^2}
\approx 126~\text{N},
\end{equation}
about 13~kg --- try it with an acrylic strip. Fixed--free drops it to
$\approx 32$~N, fixed--fixed raises it to $\approx 505$~N. The stress at
buckling ($\approx 0.63$~MPa) is far below yield, so the elastic analysis
governs; stocky columns yield first, and post-buckling behavior is beyond
this linearised model.

\vfill
{\color{celline}\hrule height 0.6pt}
\smallskip
{\footnotesize\textcolor{celdark}{Part of the multimodal mechanics series
(analytical \textperiodcentered\ simulated \textperiodcentered\ experimental)
\,\textperiodcentered\, AIM Group, Imperial College London.}}

\end{document}
